\chapter{Your First Program}\label{ch:first}

\begin{goals}
\begin{itemize}
  \item What a program is and how a computer runs it
  \item Getting to know the Salam online editor and running your first program
  \item Printing text and numbers on the screen with \code{println} and \code{print}
  \item Writing comments inside a program
  \item Reading an error message and fixing the simplest mistakes
\end{itemize}
\end{goals}

\section{What is a program?}

Suppose you want to teach a friend who has never made tea how to make it.
You cannot just say ``make a good cup of tea''. You have to go step by step:
fill the kettle with water, put it on the stove, wait until it boils, put the
tea leaves in the pot, and so on to the end. A computer is exactly this kind
of friend: very fast, very precise, and not very bright. It will do anything
you tell it, but only if you say precisely what to do.

A \textbf{program} is this list of step-by-step instructions, and
\textbf{programming} is the art of writing it. The language we write the
instructions in is called a \textbf{programming language}. Unlike a human
language, a programming language has very few words and very strict rules,
but in return it has no ambiguity at all. The language of this book is
\textbf{Salam}.

\begin{note}
The text you write when you program is called \textbf{code}, or
\textbf{source code}. The computer reads this text, carries out its
instructions one by one, and shows the result. Doing this is called
\textbf{running} the program.
\end{note}

\section{The Salam online editor}

You do not need to install or set up anything to write and run the programs
in this book. Open your browser and go to this address:
\begin{center}
\url{https://editor.salamlang.ir}
\end{center}
The page you see is called the Salam playground, and it has three main
parts:
\begin{enumerate}
  \item \textbf{The top bar.} Here you choose the language. Pick
        \textbf{English} so that both the page and the editor are set up
        for English code. The same bar offers two modes, \textbf{Program}
        and \textbf{Layout}; throughout this book we work in
        \textbf{Program} mode. The \textbf{Run} button is here too.
  \item \textbf{The editor.} The window where you write your code. Line
        numbers appear beside it, and the important words of the language
        are shown in colour.
  \item \textbf{The output.} The window where the result of running your
        program appears. If the program has a mistake, the error message
        appears here as well.
\end{enumerate}

\begin{tip}
The top bar of the editor also holds a list of sample programs. It is worth
opening one of them right now and pressing Run, just to see that everything
works. Do not worry if you cannot make sense of the code yet. By the end of
this book you will understand all of it.
\end{tip}

\section{Hello, World!}

There is an old tradition among programmers: the first program in any new
language is one that writes the sentence ``Hello, World!'' on the screen. We
will keep that tradition. Clear the editor and type these three lines into
it:

\program{Hello, World!}
\begin{salamcode}
func main:
    println "Hello, World!"
end
\end{salamcode}

Now press \textbf{Run}. (If the \textbf{Auto} switch in the top bar is on,
which it is by default, the editor runs the program by itself as you type.)
In the output window you will see this:

\begin{salamout}
Hello, World!
\end{salamout}

Congratulations! You have just written and run your first program. Now let
us see what these three lines say.

\subsection{func main: the starting point}

The first line, \code{func main:}, tells Salam: ``the program starts
running here.'' Every Salam program must have exactly one \code{func main}.
The word \code{end} on the last line marks where it finishes. For now, think
of these two lines as the covers of a notebook: whatever we write, we write
between them. In Chapter~\ref{ch:functions} we will see that the word
\code{func} has a good deal more to offer.

Look at the colon at the end of the first line. In Salam, wherever a block (a
group of instructions that belong together) begins, the line ends with
\code{:} and the block is closed with \code{end}. You will see this pattern
again and again.

\subsection{Indentation: the visual order of a program}

Notice that the middle line is set a little further in than the other two.
This is called \textbf{indentation}, and it shows that the line is
\emph{inside} \code{main}. The custom is to indent everything inside a block
by four spaces. Salam does not need these spaces to run the program, and the
program runs even if you leave them out; but a program without proper
indentation is like a book without paragraphs: it can be read, but with
difficulty.

\subsection{println: writing on the screen}

The middle line is where the work is done. The instruction \kwd{println}
writes whatever follows it in the output window and then moves to the next
line. The text we want printed goes between two quotation marks \code{"}.
Text between quotation marks is called a \textbf{string}; we will have a lot
to do with strings in Chapter~\ref{ch:strings}. The quotation marks themselves
are not printed. They only tell Salam where the text starts and where it
ends.

\begin{tip}
The quotation mark Salam expects is the straight double quote \code{"},
which on most keyboards is typed by holding Shift and pressing the \code{'}
key. If you paste code from a word processor or a chat program, check that
its ``smart'' curly quotes have not replaced the straight ones; Salam does not
treat them as quotation marks. Inside a string, however, you can write any
character you like, including curly quotes: \code{println "She said:
\textquotedblleft hi\textquotedblright"} is perfectly fine.
\end{tip}

\section{Several instructions in a row}

A program rarely has just one instruction. We write the instructions one
under another, and Salam runs them in that order, from top to bottom:

\program{Several lines of output}
\begin{salamcode}
func main:
    println "Welcome to the Salam programming book."
    println "This is the second line."
    println "And this is the third."
end
\end{salamcode}

\begin{salamout}
Welcome to the Salam programming book.
This is the second line.
And this is the third.
\end{salamout}

Each \code{println} produces a new line of output. Order matters: if you swap
two instructions, the two lines in the output swap as well. Try it right now.

\section{Printing numbers and doing arithmetic}

\code{println} is not only for text; it prints numbers too. Numbers are
written without quotation marks:

\begin{salamcode}
func main:
    println 1404
    println 2 + 3
    println 10 * 4
end
\end{salamcode}

\begin{salamout}
1404
5
40
\end{salamout}

Look at the second line: we wrote \code{2 + 3}, but Salam printed \code{5}.
Salam did the arithmetic before printing. The sign \code{*} is what
programmers use for multiplication, because there is no \texttimes{} key on
the keyboard. Chapter~\ref{ch:operators} covers arithmetic in full.

Now see what happens if we put the same \code{2 + 3} inside quotation marks:

\begin{salamcode}
func main:
    println "2 + 3"
    println 2 + 3
end
\end{salamcode}

\begin{salamout}
2 + 3
5
\end{salamout}

Whatever is inside quotation marks is \emph{text}, and Salam does not care
what it says; it prints it just as it is. Outside the quotation marks,
\code{2 + 3} is a calculation. This difference between ``text'' and
``number'' is one of the most important things you will learn in this book.

\section{Printing several things on one line}

To print both text and a number on the same line, we separate them with a
comma \code{,}. Salam prints them one after another with a single space
between them:

\program{Text and a number on one line}
\begin{salamcode}
func main:
    println "The sum of 2 and 3 is", 2 + 3
    println "Sara's age:", 17, "years"
end
\end{salamcode}

\begin{salamout}
The sum of 2 and 3 is 5
Sara's age: 17 years
\end{salamout}

\section{print: writing without moving to the next line}

The instruction \kwd{print} is just like \code{println}, with one
difference: after writing, it does not move to the next line. The next
instruction carries on from the same place:

\program{The difference between print and println}
\begin{salamcode}
func main:
    print "a"
    print "b"
    print "c"
    println ""
    println "d"
    println "e"
end
\end{salamcode}

\begin{salamout}
abc
d
e
\end{salamout}

The three \code{print} instructions all wrote on one line. Then
\code{println ""} printed an empty string, which means it wrote nothing but
moved to the next line. \code{print} becomes useful in later chapters, when
we want to build a line piece by piece. Until then we will mostly use
\code{println}.

\begin{tip}
\code{println ""} on its own produces an empty line in the output, and it is
a handy way to put some space between the parts of your output. Salam
insists on the pair of quotation marks: a bare \code{println} with nothing
after it is an error.
\end{tip}

\section{Writing comments in a program}

Sometimes we want to leave a note in the middle of the code, for ourselves or
for others: what this line does, why it is written this way, or what is not
finished yet. Such notes are called \textbf{comments}. Salam ignores
everything that follows two slashes \code{//} up to the end of that line:

\program{Comments in a program}
\begin{salamcode}
// This program prints a welcome message
func main:
    println "Welcome"    // print the main message
    // println "This line does not run"
    println "Have a nice day"
end
\end{salamcode}

\begin{salamout}
Welcome
Have a nice day
\end{salamout}

Look at the fourth line: although it is a \code{println} instruction, it did
not run, because we put \code{//} in front of it. Programmers use this trick
to switch a line off temporarily without deleting it.

If a comment runs over several lines, you can write it between \code{/*} and
\code{*/}:

\begin{salamcode}
/*
   This program was written by Sara.
   Its job: print a short poem by Saadi.
*/
func main:
    println "Human beings are members of a whole,"
    println "In creation of one essence and soul."
end
\end{salamcode}

\begin{salamout}
Human beings are members of a whole,
In creation of one essence and soul.
\end{salamout}

\begin{tip}
A good comment says something that is not obvious from the code itself: it
says \emph{why}, not \emph{what}. Writing \code{// print the message} above
\code{println "message"} helps nobody; writing \code{// this figure comes
from the 2026 price list} is valuable. In this book we comment more than
usual, because we are teaching.
\end{tip}

\section{When we make mistakes}

Everyone makes mistakes; experienced programmers simply find theirs faster.
Let us make a mistake on purpose and see what Salam says. We leave out the
closing quotation mark:

\begin{salamcode}
func main:
    println "Hello, World!
end
\end{salamcode}

When you press Run, instead of output you see a message like this:

\begin{salamout}
error: unterminated string literal
 --> main.salam:2:13
\end{salamout}

The error message tells us two important things: \textbf{what} went wrong
(an ``unterminated string literal'', that is, a string with no closing
quotation mark) and \textbf{where} (line 2, column 13; exactly where the
string begins). \code{main.salam} is the name the editor gives the file that
holds your program. Below the message, the editor also shows that line of
the program, with a \code{\textasciicircum} mark pointing at the place of
the mistake. Get into the habit of reading error messages carefully. They
are not your enemy, they are your guide. Appendix~\ref{app:errors} lists the
most common errors and what they mean.

A few very common mistakes in these first steps:
\begin{itemize}
  \item Forgetting the \code{end} at the end of \code{main}.
  \item Forgetting the colon after \code{func main}.
  \item Forgetting the comma between the things printed by one
        \code{println}.
  \item Misspelling an instruction, for example \code{printn} instead of
        \code{println}.
\end{itemize}

\begin{deep}
How does a computer that understands only zeros and ones run the text we
write? The answer is a program called a \textbf{compiler}. The Salam compiler
reads your code, first checks that it follows the rules of the language
(this is where you see the error messages), and then turns it into
instructions the computer's processor understands directly. In the online
editor the same compiler, running inside your browser, checks the code, and
then a program called an \textbf{interpreter} carries the instructions out
one at a time; the result is the same, only the road to it is different.
It is worth knowing that the Salam compiler is itself written in Salam.
\end{deep}

\section{A more complete program}

Let us put everything we learned in this chapter together in one program: a
small identity card that introduces you.

\program{An identity card}
\begin{salamcode}
// An identity card
func main:
    println "================="
    println "Name: Sara Ahmadi"
    println "City: Kashan"
    println "Age:", 2026 - 2009, "years"
    println "Interest: programming"
    println "================="
end
\end{salamcode}

\begin{salamout}
=================
Name: Sara Ahmadi
City: Kashan
Age: 17 years
Interest: programming
=================
\end{salamout}

Look at the ``Age'' line: instead of writing the number 17, we wrote the
difference of two years and let Salam work it out. This is a good habit.
Wherever a number comes from a calculation, let the computer do the
arithmetic, not you.

\begin{summary}
\begin{itemize}
  \item A program is a list of instructions that the computer carries out
        in order, from top to bottom.
  \item Every Salam program begins with \code{func main:} and finishes
        with \code{end}. The instructions go between these two lines,
        indented.
  \item \code{println} writes something and moves to the next line;
        \code{print} writes but stays on the same line.
  \item Text goes between \code{"} marks, numbers do not. Several things
        are separated by \code{,}.
  \item Anything after \code{//} is a comment and does not run.
  \item An error message says what went wrong and on which line. Read it
        carefully.
\end{itemize}
\end{summary}

\section*{Exercises}
\markright{Exercises}

\exercise Write a program that prints your first name, your last name and
your city, each on a separate line.

\exercise Write a program that prints the sentence ``I am a programmer'' on
one line using three \code{print} instructions and one \code{println}.

\exercise Write a program that does these calculations and prints each
result on its own line, with a short description: the sum of 125 and 375;
the product of 24 and 7; the number of hours in a week.

\exercise The program below has three mistakes. Find them without running
it, then run it to check your answer:
\begin{salamcode}
func main
    pritnln "Hello"
    print "Goodbye
end
\end{salamcode}

\exercise Write a program that prints a four-line poem or proverb of your
choice, and in a multi-line comment at the top of the program, write the
name of the poet or the source.
